\begin{frame}{``하나를 고르는 문제''에서 ``분포의 형태''로}
  \begin{itemize}
    \item $n$-step이 \textbf{``하나의 $n$을 고르는''} 문제였다면,
      TD($\lambda$)는 \textbf{``가중치 분포 $\lambda^{n-1}$의 형태''}를
      $\lambda$ 하나로 조절하는 문제.
    \item 같은 편향-분산 트레이드오프를, \textbf{하나를 고르는 문제에서
      분포의 형태를 고르는 문제로 승화}시킨 것.
  \end{itemize}
  \vskip0.8em
  \begin{block}{직접 확인해야 할 한 줄}
    $\lambda$를 $n$으로 환산하는 근사 대응: TD($\lambda$)의
    \textbf{유효 수명(effective horizon) $\sim 1/(1-\lambda)$ 스텝}.
    $\lambda=0.8$이면 대략 $1/(1-0.8) = 5$ 스텝 전후 --- 이번 실험에서
    $n=5$가 최적, $\lambda=0.8$이 최적인 것과 \textbf{동일한 방향}.
    정확한 1:1 변환은 아니지만, ``$\lambda=0.8$이 5스텝 정도
    내다보는 것과 비슷하다''는 직관은 잡아둠.
  \end{block}
\end{frame}
